\section{Introduction}\label{sec:intro}

\subsection{Overview of the main contributions}

The results of \cite{noferini2016numerical} and \cite{graf2026numerical} say that every known
algebraic rootfinder is exponentially unstable in the worst case: there exist systems for which
the intermediate subproblem (an eigenvalue problem or a univariate polynomial) is worse
conditioned than the original root by a factor $\sigma^{-(d-1)}$, or even $\sigma^{-2^d+2}$ for
elimination. These statements rest on carefully designed examples. This paper asks what happens
for typical systems. We model ``typical'' by i.i.d.\ complex Gaussian coefficients with fixed
degree $n$ in each of the $d$ variables, the same degree model chosen in
\cite{noferini2016numerical}, and we study the expected value of the logarithm of the
conditioning ratio, averaged over all roots of the system. Probabilistic analysis of this
kind has long been central to the other main family of polynomial system solvers, homotopy
continuation methods: the average-case programme initiated by Shub and Smale
\cite{shub1993complexity}, and systematized in the monographs of Blum, Cucker, Shub, and
Smale \cite{blum1998complexity} and of B\"urgisser and Cucker \cite{burgisser2013condition},
led, through the probabilistic solution of Smale's
seventeenth problem by Beltr\'an and Pardo \cite{beltran2009smale} and the deterministic
near-solution of B\"urgisser and Cucker \cite{burgisser2011smale}, to Lairez's
deterministic average polynomial-time algorithm \cite{lairez2017deterministic}. The present
paper brings the same average-case viewpoint to the algebraic rootfinders.

The first outcome is structural. For every method in the family, the ratio at a root factorizes
into a \emph{local} part, depending only on the root and on the system near it, and a
\emph{structural} part, encoding how the method couples distinct roots. The local part is the
same for all methods, and we compute it exactly: it equals $(d+1)\log n$ up to an additive
$O(d)$. It should be read as a change of units between ``conditioning of the root'' and
``conditioning of the subproblem'', not as an instability; a scale-free version of the
comparison removes it entirely (Remark~\ref{rem:tworeadings}).

The second outcome is a hierarchy. The structural parts fall into three classes. For the
hidden-variable resultant methods in their canonical Cayley and Sylvester forms there is none:
an exact identity of \cite{noferini2016numerical} makes the ratio a purely local quantity, and
for Cayley we obtain a complete and unconditional average-case theorem, verified numerically,
with the conclusion that the exponential worst-case instability of \cite{noferini2016numerical}
is a tail event, not the typical behaviour; in the relative metric the method is, on average,
neutral with respect to the root conditioning. Multiparameter eigenvalue problems are also
purely local. For Gr\"obner elimination and the rational univariate representation, the
structural part is a logarithmic energy of pairs of roots, which requires only the two-point
statistics of the root set. For normal form and Macaulay methods it is the eigenvector
conditioning of the multiplication matrix, a quantity involving all roots at once. We prove
that this hardest term is finite for every $d$, and the numerical protocol, executed at
$d=1,2,3$, shows that it does not follow the random-matrix calibration at $d=1$ and $2$, where
it appears bounded and the multiplication matrix is far more rigid than a Ginibre matrix; at
$d=3$ the measured overlaps run parallel to the Ginibre control, a constant below it, up to
$N=1296$: the third dimension appears to sit in the growth regime at accessible sizes. If so,
the average loss of the normal-form method is fully described by the local terms; the
pairwise-class terms are computed at leading order: one is an outer-root effect, the other a
closed-form logarithmic energy. Section~\ref{sec:relative} converts everything to the relative
metric, the one in which conclusions should be drawn, and the family separates cleanly there:
only the rational univariate representation pays an exponential average penalty, with a
closed-form constant. The message is that the exponential instability of the worst case is
produced by atypical configurations, chiefly nearly colliding roots; for random systems the
expected loss is modest, in agreement with our earlier numerical evidence.

A byproduct of independent interest is what we call the outer-root phenomenon. For Kac-type
random systems, logarithmic averages over the roots are not dominated by the roots near the
unit torus. About half of the roots lie outside, at rescaled distances with a heavy $1/u$ tail,
and the local scales of the system at such roots are exponentially large. Their total
contribution to every ``average of a logarithm'' is of the same order as that of the bulk.
Several natural heuristics fail at leading order for exactly this reason. The Kostlan ensemble
exhibits the same effect in extreme form: its root measure does not concentrate at all, the
affine average ratio grows linearly in $n$, and the projectively invariant normalization is
identified as the only sensible yardstick there. The total-degree model, by contrast,
reproduces the separate-degree coefficients exactly, through a different mechanism: its heavy
tail is carried by roots receding radially in all coordinates at once.

All local components are established rigorously, as is the finiteness of every structural term,
and the overlap term is proved bounded at $d=1$ and $d=2$; what remains open is the sharp
growth of the full-configuration terms, together with the rigor items collected in
Section~\ref{sec:open}.

\subsection{The quantity of interest}\label{sec:quantity}

Let $P=(p_1,\dots,p_d):\C^d\to\C^d$ be a square polynomial system with $N$ simple roots
$z_1,\dots,z_N$. An algebraic rootfinder computes the roots through an intermediate problem:
normal-form methods diagonalize a multiplication matrix on the quotient ring; hidden-variable
resultant methods solve $\det R(x_1)=0$ for a resultant matrix polynomial; Gr\"obner
elimination and the rational univariate representation produce a univariate polynomial;
multiparameter eigenvalue methods and the Macaulay resultant produce generalized eigenvalue
problems. Each method M has a condition number $\kappa_{\mathrm{M}}(x^*)$ for the subproblem it
constructs at a root, defined in \cite{graf2026numerical,noferini2016numerical}. The intrinsic
condition number of the root is $\kroot(x^*):=\|J(x^*)^{-1}\|_2 = 1/\sigma_{\min}(J(x^*))$,
with $J$ the Jacobian. The object of study is the average logarithmic conditioning ratio
\begin{equation}\label{eq:ratio}
  \Rcal_n^{\mathrm{M}} := \frac1N\, \E\Bigl[\sum_{j=1}^N \log
  \frac{\kappa_{\mathrm{M}}(z_j)}{\kroot(z_j)}\Bigr],
\end{equation}
where the expectation is over the random coefficients (Section~\ref{sec:model}) and the sum
runs over all roots. We use $\E[\log]$ rather than $\log\E$ because it measures the typical
number of digits lost, it is additive under the factorizations below, and it is robust to the
heavy tails that condition numbers have.

All the methods above are treated, and the plan rests on one structural fact, proved case by
case below: for every method, the ratio in \eqref{eq:ratio} factors at each root into a
\emph{local} part, a function of the root and of the system near it, and a \emph{structural}
part, encoding how the method couples distinct roots. The local part turns out to be the same
for all methods, so it is computed once, in
Sections~\ref{sec:jacobian}--\ref{sec:local}; the methods differ only in their structural
parts, which are classified and analysed in Sections~\ref{sec:structural}
and~\ref{sec:family}. The common machinery is developed first, and among the structural parts
the normal form's is treated before the others, for two reasons. First, its condition number,
by \cite[Thm.~6.2]{graf2026numerical},
\begin{equation}\label{eq:keig}
  \keig(x^*) = \frac{\|[\det Q]_\B\|_2\, \|\B(x^*)\|_2}{|\det J(x^*)|}\,(1+|\lambda|),
\end{equation}
where $\B$ is a monomial basis of the quotient ring $R/I$, $Q$ is any matrix polynomial with
$P(x)=Q(x)(x-x^*)$, and $\lambda=g(x^*)$ is the eigenvalue of the multiplication matrix $M_g$
associated with the root ($g$ a fixed generic polynomial such as $g=x_1$), leads to the first
of the two exact reductions on which everything rests. Second, its structural part is the only
genuinely global one, the eigenvector conditioning of $M_g$, so it is the case that needs the
full machinery; the other methods are then treated on the same footing, at little extra cost,
and some, like the Cayley resultant, end up completely solved. We write
$\Rcal_n := \Rcal_n^{\mathrm{eig}}$ for the normal-form instance.

To avoid any ambiguity about which condition numbers enter \eqref{eq:ratio}: for the normal
form we take \eqref{eq:keig}, the exact formula of \cite[Thm.~6.2]{graf2026numerical}, as the
definition; for the hidden-variable resultants we use the eigenvalue condition number of the
resultant matrix polynomial as defined in \cite[Def.~2.10]{noferini2016numerical}, for which
\cite[Lem.~2.11]{noferini2016numerical} is an equality; for Gr\"obner elimination and the
rational univariate representation we use the conditioning $1/|g'|$ of the univariate
subproblem, as in \cite{graf2026numerical}. All of these are actual condition numbers, not
bounds. The single exception is the Macaulay method, for which
\cite[Thm.~7.1]{graf2026numerical} provides only a lower bound; the corresponding statements
in Section~\ref{sec:macaulay} are lower bounds and are marked as such.

\subsection{Root averages carry a Jacobian weight}

One identity organizes the whole computation. For any measurable quantity $\varphi(X;P)$
depending on a point and on the system,
\begin{equation}\label{eq:counting}
  \E\Bigl[\sum_{z:\,P(z)=0} \varphi(z;P)\Bigr]
  = \int_{\C^d} \E\bigl[\varphi(X;P)\,|\det J(X)|^2 \bigm| P(X)=0\bigr]\, p_{P(X)}(0)\,dX,
\end{equation}
where $p_{P(X)}(0)$ is the value at the origin of the probability density of the Gaussian
vector $P(X)\in\C^d$. This is a change of variables (the coarea formula), not a probabilistic
subtlety; the factor $|\det J|^2$ is the real Jacobian determinant of the map $\C^d\to\C^d$,
and the square comes from holomorphy. Identity \eqref{eq:counting} is, in other words, the
Kac--Rice formula for the zero set of the Gaussian field $P$, in the weighted (or marked)
form in which the summand $\varphi$ may depend on the whole field, specialized to the
holomorphic case: the Kac--Rice weight $|\det DF|$ of a real field $F$ becomes $|\det J|^2$,
which is sign-definite and polynomial in the field and its derivatives, so none of the
absolute-value smoothness issues of the real theory arise. We refer to
Aza\"is--Wschebor \cite[Ch.~6]{azais2009level} for the general theory, and to
Berzin--Latour--Le\'on \cite{berzin2022kacrice} for a contemporary overview of the formula's
variants and applications. The practical content
is that ``conditioning on a root at $X$'' is not the naive conditioning on $\{P(X)=0\}$: the
correct conditional law -- in point-process language, the Palm distribution of the root set
-- carries the weight $|\det J(X)|^2$, which down-weights badly conditioned roots.
Forgetting the weight gives wrong answers; keeping it turns out to simplify the
computations, because the weighted Gaussian ensembles that appear are again classical
(Section~\ref{sec:weight}). The same mechanism underlies the moment computations of Shub and
Smale \cite{shub1993complexity} and, at $d=1$, the root-density computations of Edelman and
Kostlan \cite{edelman1995howmany}. Hypotheses and an extension to unbounded $\varphi$ are
checked in
Section~\ref{sec:counting}, where the identity is deduced from the Rice formulas of
\cite[Thms.~6.2 and~6.4]{azais2009level}.

\subsection{Results}\label{sec:results}

Items (a)--(c) below are common to all methods; (d) and (e) concern the structural part of the
normal form; (f) the other methods.

\begin{enumerate}[(a)]
  \item \emph{Reduction I} (Section~\ref{sec:reduction1}, exact). In $R/I$ one has
    $[\det Q] = \det J(z_j)\,[\ell_j]$, with $\ell_j$ the Lagrange idempotent at $z_j$. Hence
    $\keig(z_j) = \Lambda_j\,(1+|\lambda_j|)$, where $\Lambda_j\ge1$ is exactly the Wilkinson
    condition number of the eigenvalue $\lambda_j$ of $M_g$, and
    \[
      \frac{\keig(z_j)}{\kroot(z_j)}
      = \sigma_{\min}(J(z_j))\cdot\Lambda_j\cdot(1+|\lambda_j|).
    \]
    The reduction of $\det Q$ modulo $I$, which had appeared to be the main obstacle,
    disappears from the problem.
  \item \emph{Reduction II} (Sections~\ref{sec:jacobian}--\ref{sec:weight}, exact).
    Conditionally on $P(X)=0$, the Jacobian is $G\,\Sigma(X)^{1/2}$ with $G$ a $d\times d$
    matrix of i.i.d.\ standard complex Gaussians and $\Sigma(X)$ explicit and diagonal in our
    model. The weight $|\det G|^2$ maps the singular values of $G$ to those of a
    $(d+1)\times d$ Gaussian matrix. Every local root average is thereby an explicit integral.
  \item \emph{Local terms} (Section~\ref{sec:local}). With
    $A_n:=\frac1N\E\sum_j\log\sigma_{\min}(J(z_j))$ and
    $C_n:=\frac1N\E\sum_j\log(1+|\lambda_j|)$,
    \[
      A_n = (d+1)\log n + O(d), \qquad C_n = \log2 + o(1).
    \]
    The computation requires the exact law of the root radii at scale $1/n$
    (Proposition~\ref{prop:rootscale}) and reveals the outer-root phenomenon
    (Section~\ref{sec:outer}).
  \item \emph{The structural term} (Section~\ref{sec:structural}). Everything nonlocal sits in
    $\Theta_n := \frac1N\E\sum_j\log\Lambda_j\ge0$. We give exact identities, prove
    $\Theta_n<\infty$ for every $d$ (Proposition~\ref{prop:finiteness}), identify the
    mechanism of its tails, prove $\Theta_n=O(1)$ at $d=1$ and $d=2$
    (Theorem~\ref{thm:overlap}), and execute the numerical test at $d=1,2,3$: at $d=3$ the
    overlaps run parallel to the Ginibre control, leaving the regime open there.
  \item \emph{Main theorem} (Section~\ref{sec:theorem}, conditional):
    \[
      \Rcal_n = (d+1)\log n + \Theta_n + O(d), \qquad 0\le\Theta_n<\infty,
    \]
    and, by Theorem~\ref{thm:overlap}, $\Rcal_n = (d+1)\log n + O(1)$ unconditionally at
    $d=1,2$; the regime at $d\ge3$ is open.
  \item \emph{The full family} (Section~\ref{sec:family}). The other methods of
    \cite{graf2026numerical} obey the same skeleton with structural terms classified as: none
    (hidden-variable resultants in their canonical Cayley and Sylvester forms, by an exact
    identity from \cite{noferini2016numerical}; multiparameter eigenproblems), pair energies
    (Gr\"obner elimination, rational univariate representation), full-configuration overlaps
    (Macaulay). For the Cayley method this yields an unconditional theorem: the average ratio
    equals $-\frac{d(d-1)}{2}\log n + O_d(1)$ (Theorem~\ref{thm:cayley}), confirmed
    numerically at $d=2$ down to the finite-$n$ constants (Section~\ref{sec:verification});
    in the absolute metric of \cite{noferini2016numerical} the eigenproblem is typically
    \emph{better} conditioned than the root, a units effect that disappears in the relative
    metric, where the average ratio is neutral (Section~\ref{sec:relativeCayley}).
\end{enumerate}

\subsection{Notation}

$\|\cdot\|_2$ denotes the Euclidean norm of a vector and the spectral norm of a matrix,
$\|\cdot\|_F$ the Frobenius norm, and $\sigma_1\ge\dots\ge\sigma_{\min}$ the singular values.
$G$ always denotes a $d\times d$ matrix with i.i.d.\ $\CN(0,1)$ entries, and
$\EW[\cdot] := \E[\,\cdot\,|\det G|^2]/d!$ the associated weighted expectation
(Section~\ref{sec:weight}). The roots of the system are $z_1,\dots,z_N$ with $N=d!\,n^d$;
$\Scal$ is a monomial staircase, $\B$ the corresponding basis of $R/I$,
$V=(z_j^\beta)_{\beta\in\Scal,\,j}$ the evaluation matrix with columns $v_j$, $\ell_j$ the
Lagrange idempotents with coordinate vectors $V^{-T}e_j$, and
$\Lambda_j := \|V^{-T}e_j\|_2\,\|Ve_j\|_2$. $\psi=\Gamma'/\Gamma$ is the digamma function,
$H_d=\sum_{k=1}^d 1/k$, $\gamma$ Euler's constant, and $k_m(r):=\sum_{j=0}^m r^j$. Following
\cite{noferini2016numerical}, $\Omega$ denotes the closed unit polydisk
$\bar{\mathbb{D}}^d$, the reference domain of the rootfinding problem, so that
$\sup_\Omega|p_i|$ is the sup-norm of $p_i$ there (equal to the sup on the torus
$\mathbb{T}^d$ by the maximum principle). Asymptotic notation $O_d(\cdot)$ indicates
constants depending on $d$ only; all logarithms are natural.
